
\subsubsection{Architecture Detection Methods}

Prior methods for neural network architecture detection and model fingerprinting primarily rely on activation patterns or forward passes. These approaches generate activation patterns by passing test data through the network, then analyze the patterns for architecture identification or tampering detection. While such methods can achieve high accuracy, they require computational resources for inference and access to appropriate test inputs. The approach in this study differs by requiring only access to weight files, with no forward passes or test data.

\subsubsection{Weight Distribution Analysis}

Research on neural network weight distributions spans compression and transfer learning applications. Pruning methods analyze weight magnitudes to identify parameters that can be removed with minimal accuracy loss. Quantization work studies weight value distributions to determine optimal bit-width assignments for model compression. Transfer learning research measures weight similarity metrics to predict how well a model pretrained on one task will transfer to another task. These studies treat weight distributions as optimization targets rather than as information sources about architectural properties.

\subsubsection{Training Dynamics Research}

Theoretical analyses of deep network training examine how gradients propagate through multiple layers during backpropagation. This research has studied phenomena including vanishing and exploding gradients, and has analyzed how architectural innovations like residual connections affect gradient flow. Such work suggests that training through different network depths might create characteristic weight patterns. The random initialization tests in this study directly examine this hypothesis by comparing pretrained and untrained models with identical architectures.

\subsubsection{Architectural Innovations}

Modern convolutional neural network architectures incorporate specific structural patterns. ResNet introduced residual connections enabling training of very deep networks through skip connections. Bottleneck layers using 1×1 convolutions reduce parameter counts while maintaining representational capacity. DenseNet introduced dense connections where each layer connects to all subsequent layers within a block. These architectural components create structural patterns that may be detectable in weight statistics.
